% =====================================================================
%  BAB I PENDAHULUAN
% =====================================================================

\chapter{PENDAHULUAN}

\section{Latar Belakang}

Pencatatan berat barang di gudang dan pabrik masih banyak dilakukan secara
manual. Operator membaca angka pada tampilan timbangan digital, kemudian
menuliskannya ke buku catatan atau lembar kerja. Cara tersebut memiliki dua
kelemahan yang melekat. Pertama, tidak tersedia pembanding otomatis, sehingga
kesalahan catat tidak dapat dideteksi pada saat kejadian. Kedua, tidak
tertinggal jejak audit, sehingga sengketa jumlah barang sulit ditelusuri.
Kebutuhan akan perekaman otomatis telah dijawab di literatur, antara lain
melalui sistem pencatat data waktu nyata berbasis \textit{Convolutional
Recurrent Neural Network} dan \textit{Connectionist Temporal Classification}
(CRNN-CTC) untuk tampilan tujuh segmen \parencite{effendi2025unjuk}. Karena
tampilan timbangan sudah menyajikan angka secara digital, pembacaan sebenarnya
dapat diotomasi menggunakan kamera dan pengenalan karakter optis tanpa perlu
mengganti perangkat timbangan yang telah terpasang.

Persoalannya, mesin pengenalan karakter yang dirancang untuk teks dokumen tidak
bekerja dengan baik pada tampilan tujuh segmen. \textcite{low2024detecting}
melaporkan bahwa pengenalan pada kumpulan citra tujuh segmen tetap bermasalah,
terutama pada kondisi pantulan cahaya. Penyebabnya bersifat struktural:
karakter pada tampilan tujuh segmen bukan huruf, melainkan kombinasi segmen yang
menyala. Angka ``1'' hanya berupa dua garis tegak, sedangkan angka ``8''
menyalakan ketujuh segmen. Jarak antardigit sangat rapat, dan titik desimal
berukuran hanya beberapa piksel sehingga sering tidak terbaca sebagai karakter
tersendiri. Persoalan desimal inilah yang menyebabkan
\textcite{gushima2023automatic} memisahkan model digit dari model titik desimal.

Sejumlah pendekatan khusus kemudian dikembangkan. \textcite{laroca2019convolutional}
membangun metode berbasis jaringan konvolusional dan merilis kumpulan data
UFPR-AMR berisi 2.000 citra meter energi beserta 10.000 anotasi tingkat digit.
\textcite{bonilla2023adapting} mengadaptasi YOLO untuk mengenali angka riil
termasuk titik desimal. \textcite{yenneti2023lightweight} mengusulkan kepala
inferensi digit yang mencapai 97,8 persen pada kerangka tunggal dengan kecepatan
30 bingkai per detik, disertai penyaringan antarbingkai.
\textcite{wei2023novel} melengkapinya dengan model sekuens bermemori pendek
untuk panjang bacaan yang bervariasi. Di sisi lain,
\textcite{jiang2026method} justru menunjukkan bahwa pendekatan klasik berbasis
binarisasi adaptif, proyeksi, dan pencocokan Hamming mencapai 99,7 persen dalam
20 milidetik pada timbangan elektronik. Temuan tersebut menandakan bahwa
pembelajaran mendalam tidak serta-merta mengungguli metode klasik pada domain
yang geometrinya sangat teratur.

Pada periode berikutnya perhatian bergeser ke model yang lebih besar dan lebih
umum. \textcite{ega2026leveraging} membandingkan model bahasa-visi dengan lima
mesin pengenalan karakter pada perangkat ukur digital dan melaporkan bahwa
Gemini mencapai \textit{full-string recognition} sebesar 93,85 persen.
Angka tersebut tampak bertentangan dengan temuan \textcite{lin2026measurebench}
yang melaporkan model bahasa-visi terbaik hanya mencapai 30,2 persen pada data
nyata. Keduanya tidak bertentangan: pengujian pertama menilai transkripsi digit
yang telah tertulis, sedangkan pengujian kedua menuntut inferensi skala dan
penalaran spasial. Perbedaan itu menempatkan model bahasa-visi bukan sebagai
pesaing akurasi pada domain transkripsi digit, melainkan sebagai pesaing biaya.

Pada Juni 2026 dirilis PP-OCRv6, keluarga model pengenalan karakter ringan
berukuran 1,5 hingga 34,5 juta parameter \parencite{zhang2026ppocrv6}. Model
tersebut memakai tulang punggung PPLCNetV4, leher pengenalan
\textit{EncoderWithLightSVTR}, serta dekoder gabungan CTC dan NRTR. Pada tolok
ukur internal pengembangnya, tingkat \textit{medium} mencapai 83,2 persen
akurasi pengenalan rata-rata dari 15 kategori dan secara eksplisit menyebut
tampilan digital sebagai salah satu sasaran khusus
\parencite{zhang2026ppocrv6}. Kehadirannya melanjutkan kecenderungan yang telah
terlihat pada PP-OCRv5, yaitu model berukuran 5 juta parameter yang mampu
menyaingi model bahasa-visi berukuran miliaran parameter pada tugas serupa
\parencite{zhang2026ppocrv5}.

Seluruh pekerjaan di atas berhenti pada titik yang sama: melaporkan akurasi pada
kumpulan data yang disusun sendiri. Angka yang beredar tidak dapat dibandingkan
satu sama lain karena dihasilkan dari data, pembagian latih-uji, dan definisi
``benar'' yang berbeda. Persoalan yang lebih mendasar terletak pada apa yang
dilaporkan. Semua sumber tersebut mengukur akurasi, tetapi tidak mengukur
kesalahan yang paling mahal dalam pencatatan, yaitu bacaan yang salah namun
lolos tanpa peringatan. Pada konteks penimbangan, satu bacaan salah lebih mahal
dampaknya daripada sepuluh bacaan yang ditolak, karena bacaan salah masuk ke
catatan dan baru diketahui setelah transaksi selesai. Sistem nyata memang telah
memakai ambang keyakinan, tetapi ambang tersebut dipilih tanpa divalidasi
terhadap data, padahal nilai keyakinan model bukan bukti kebenaran.
Sepanjang penelusuran penulis terhadap 24 sumber terbitan 2019--2026, tidak
ditemukan penelitian yang memakai konteks di luar citra untuk memverifikasi
hasil pembacaan tampilan tujuh segmen.

Berdasarkan uraian tersebut, penelitian ini mengambil tiga hal yang belum
dikerjakan sumber sebelumnya. Pertama, menguji kemampuan PP-OCRv6 dalam membaca
tampilan tujuh segmen tanpa adaptasi dan menentukan adaptasi minimum yang
diperlukan \parencite{zhang2026ppocrv6}. Kedua, merancang lapisan verifikasi
kontekstual berlapis yang memanfaatkan konteks di luar citra untuk memutuskan
apakah suatu bacaan layak dipercaya. Ketiga, mengukur keandalan sistem tidak
hanya melalui akurasi, tetapi juga melalui laju kesalahan senyap,
\textit{coverage}, \textit{selective accuracy}, dan luas daerah di bawah kurva
\textit{risk--coverage}.

\section{Rumusan Masalah}

Berdasarkan latar belakang tersebut, rumusan masalah penelitian ini adalah
sebagai berikut.
\begin{poin}
  \item Sejauh mana PP-OCRv6 mampu membaca tampilan tujuh segmen tanpa
        adaptasi, dan adaptasi minimum apa yang diperlukan agar kemampuannya
        memadai?
  \item Bagaimana merancang lapisan verifikasi kontekstual yang dapat menekan
        laju kesalahan senyap pada pembacaan tampilan tujuh segmen?
  \item Berapa \textit{trade-off} antara \textit{coverage} dan akurasi pada
        bacaan yang diterima, serta bagaimana dampaknya terhadap beban koreksi
        operator?
\end{poin}

\section{Tujuan Penelitian}

Tujuan yang ingin dicapai penelitian ini adalah sebagai berikut.
\begin{poin}
  \item Mengukur kemampuan PP-OCRv6 pada tampilan tujuh segmen dalam kondisi
        tanpa adaptasi dan menentukan tingkat adaptasi minimum yang diperlukan.
  \item Merancang dan mengimplementasikan empat lapis verifikasi kontekstual,
        yaitu verifikasi struktural, fisik, temporal, dan statistik.
  \item Menghasilkan bukti kuantitatif keandalan sistem berupa laju kesalahan
        senyap, \textit{coverage}, \textit{selective accuracy}, dan luas daerah
        di bawah kurva \textit{risk--coverage}, beserta dampaknya terhadap
        tingkat koreksi operator.
\end{poin}

\section{Manfaat Penelitian}

Penelitian ini diharapkan memberikan manfaat sebagai berikut.
\begin{poin}
  \item Bagi pemilik usaha dan operator, tersedia catatan berat yang dapat
        diaudit, dan bacaan yang meragukan tidak lagi masuk ke catatan secara
        diam-diam.
  \item Bagi pengembangan ilmu, tersedia kerangka evaluasi keandalan untuk
        domain tampilan tujuh segmen, yang selama ini hanya diukur dengan
        akurasi \parencite{effendi2025unjuk,gushima2023automatic,yenneti2023lightweight}.
  \item Bagi praktik rekayasa, tersedia bukti kuantitatif mengenai biaya
        penolakan, yaitu berapa banyak bacaan yang harus dikorbankan untuk
        mencapai tingkat kepercayaan tertentu, sehingga pemilihan ambang dapat
        dilakukan berdasarkan data dan bukan perkiraan.
\end{poin}

\section{Batasan Masalah}

Agar pembahasan terarah, penelitian ini dibatasi pada hal-hal berikut.
\begin{poin}
  \item Karakter yang dikenali terbatas pada angka 0--9 dan titik desimal,
        dengan panjang 1--8 karakter.
  \item Objek penelitian adalah tampilan tujuh segmen jenis \textit{light
        emitting diode} dan \textit{liquid crystal display} pasif pada timbangan
        digital.
  \item Data uji berupa citra nyata yang tidak dipakai untuk pelatihan maupun
        penyetelan, dan dikunci sebelum eksperimen dijalankan.
  \item Data latih dihasilkan secara sintetis dengan kondisi terkontrol.
  \item Sistem dirancang berjalan mandiri di perangkat tanpa koneksi jaringan,
        dengan sinkronisasi data dilakukan saat jaringan tersedia.
  \item Lingkup sistem terbatas pada satu lini produk tanpa integrasi dengan
        sistem perencanaan sumber daya perusahaan.
\end{poin}
